\chapter{KẾT LUẬN}
\label{chap:conclusion}

\section{Tổng kết}

Báo cáo đã xây dựng và huấn luyện ViVLM-nano --- một mô hình ngôn ngữ--thị
giác tiếng Việt mà phần ngôn ngữ 101 triệu tham số được tiền huấn luyện từ
token đầu tiên trên 3,02 tỷ token FineWeb2-HQ (val loss 2,61, PPL 13,91), sau
đó hợp nhất thị giác kiểu LLaVA (SigLIP đóng băng, projector pixel-shuffle),
tinh chỉnh có giám sát hai pha (CIDEr UIT-ViIC 0,325) và tinh chỉnh học tăng
cường SCST (nâng CIDEr thêm 0,068) --- toàn bộ ở quy mô một GPU. Cùng đồ án
giữa kỳ (nhánh phân loại) và đồ án cuối kỳ
(nhánh sinh có điều kiện đơn phương thức), đề tài khép kín sơ đồ kiến trúc đa
phương thức của môn học ở phần trung tâm nhất: hợp nhất hai luồng token và
chuỗi huấn luyện pretrain--SFT--RL.

Về phương pháp làm việc, hai điểm được kiểm chứng là đáng giá: (i)~kỷ luật
kiểm thử (52 unit test ngoại tuyến, bất biến khôi phục checkpoint chính xác
từng bit) cho phép yên tâm chạy phiên huấn luyện dài; (ii)~tách bạch ba giai
đoạn bằng checkpoint độc lập giúp quy được đóng góp của từng bước vào kết quả
cuối cùng.

\section{Hạn chế}

Bốn hạn chế chính. \textbf{Quy mô}: 101M tham số và 3 tỷ token đặt trần thấp
cho kiến thức thế giới --- đặc biệt lộ ở hỏi--đáp mở; kết quả không nhằm cạnh
tranh các VLM tỷ tham số mà nhằm định lượng ``một GPU làm được đến đâu''.
\textbf{Dữ liệu SFT}: 77 nghìn mẫu là nhỏ so với hàng triệu mẫu của các VLM
công bố; hai bộ mô tả ảnh lệch miền (thể thao, đời sống) nên khả năng khái
quát ngoài miền chưa được đo. \textbf{Đánh giá}: chưa có đối chiếu trên
benchmark VLM chuẩn quốc tế (chúng chưa có bản tiếng Việt tương thích quy mô
này); BPC so với duy nhất một baseline công khai. \textbf{Bộ mã hóa thị
giác}: đóng băng hoàn toàn --- chưa khảo sát tinh chỉnh end-to-end vốn là đòn
bẩy lớn nhất ở đồ án cuối kỳ.

\section{Hướng phát triển}

Nếu tiếp tục, thứ tự ưu tiên: (i)~\emph{DPO} trên cặp ưa thích tự tạo (beam
tốt nhất/tệ nhất theo CIDEr) --- so sánh trực tiếp với SCST trong cùng ngân
sách; (ii)~\emph{KV-cache} cho giải mã và \emph{cache đặc trưng ảnh} khi đánh
giá --- hiện mỗi bước sinh chạy lại toàn chuỗi; (iii)~mở khóa dần SigLIP
(tinh chỉnh end-to-end) đối chứng với kết luận encoder-là-đòn-bẩy của đồ án
cuối kỳ; (iv)~trộn thêm dữ liệu chỉ dẫn thuần văn bản vào SFT để giữ chất
lượng ngôn ngữ khi học kỹ năng nhìn.
